% Section 6 — Cross-dataset transfer
% Page budget: ~2 pages in acmart sigconf two-column (~1800-2100 words).
% Source-of-truth: results/otel-demo/PHASE3-OTEL-SUMMARY.md,
%                  results/wol/PHASE3-WOL-SUMMARY.md,
%                  DOCS/docs8/RQ-CLOSURE-TABLE.md.
%
% Editorial discipline:
%   - The cascade is internal. The transfer claim is about the AGENT
%     and its capability-adaptive policy, not about a particular
%     retriever beating another particular retriever.
%   - All measured numbers cite a source-of-truth SUMMARY.md.
%   - The §5 OB-deep story is referenced rather than re-stated;
%     §6 is the cross-dataset replication, not a re-run of §5.

\section{Cross-Dataset Transfer}
\label{sec:cross-app}

The agent's design (\S\ref{sec:approach}) was developed against \dsob{}. The within-dataset results (\S\ref{sec:results}) show that the design's structural claims---cost savings without retrieval-quality loss, capability-driven gating, paired-delta-significant load-bearing primitives---hold on the synthetic microservices benchmark. This section asks whether they survive a change of dataset, and reports what generalises, what does not, and what the cross-dataset comparison itself reveals about the agent's failure modes.

\subsection{OpenTelemetry Demo: structurally different synthetic application}
\label{sec:cross-app:otel}

The first transfer test is to \dsotel{} (\S\ref{sec:evaluation:datasets:otel}), a polyglot synthetic application with fifteen-plus services, Kafka-mediated asynchronous communication, and a fifty-three-family scenario taxonomy more than twice the size of \dsob{}'s. The setup holds the agent's design fixed and varies only the upstream dataset: the same controller branches, the same skill registry, the same evidence-gating logic, the same on-demand tool loop, and the same statistical envelope.

Table~\ref{tab:headline-otel} reports the agent's headline metrics on the $247$-window test split. The dataset profile registers two of the cascade-derived primitives (the numeric-feature classifier and the dense retriever) plus the hybrid retriever; the four on-demand evidence-gathering tools are all available because the dataset captures the same telemetry modalities as \dsob{}.

\begin{table}[ht]
  \centering
  \caption{Headline metrics on \dsotel{} ($n=247$, $n_{\text{eval}} = 119$). $95\%$ percentile CIs from $1{,}000$-resample bootstrap at seed $42$.}
  \label{tab:headline-otel}
  \small
  \begin{tabular}{lrr}
    \toprule
    Metric & Point & $95\%$ CI \\
    \midrule
    Hit@1                     & $0.6471$ & $[0.5641, 0.7411]$ \\
    Hit@5                     & $0.7563$ & $[0.6726, 0.8333]$ \\
    MRR                       & $0.6880$ & $[0.6132, 0.7715]$ \\
    Triage accuracy           & $0.8421$ & $[0.7976, 0.8866]$ \\
    Pages-per-incident        & $1.000$  & --- \\
    Suppressions fired        & $1$      & --- \\
    Distinct plan identifiers & $4 / 6$  & --- \\
    \bottomrule
  \end{tabular}
\end{table}

The cross-dataset reading is the comparison against Table~\ref{tab:headline-ob}: Hit@5 of $0.7563$ on \dsotel{} is essentially identical to $0.7583$ on \dsob{}, with overlapping confidence intervals, despite the dataset being structurally different and the controller and skill registry being unchanged. Hit@1 is $0.04$ lower in absolute terms (a difference within the per-dataset CI width). Triage accuracy is marginally higher ($0.8421$ versus $0.8343$), again within noise. The agent's design transfers across applications without any per-application tuning.

\paragraphhead{Cost.} On \dsotel{} the agent saves $65.3\%$ of wall time, language-model tokens, and dollar-equivalent cost against the always-on cascade counterfactual---higher than \dsob{}'s $64.1\%$. The mechanism is the same: the cheap-path-confident gate fires on more windows on \dsotel{} because the dataset's triage-numeric confidence distribution is more bimodal (most windows are clearly noise or clearly an active fault, with fewer ambiguous cases that would force the gate open). The cost-savings claim survives the dataset change.

\paragraphhead{Pareto curve.} The threshold sweep over $\{0.50, \ldots, 0.95\}$ produces an even flatter curve on \dsotel{} than on \dsob{}: every threshold produces identical Hit@K \emph{and} identical cost, because the small evaluable subset (\nopagebreak $n=119$) and the bimodal cheap-path confidence distribution leave no windows in the threshold-sensitive intermediate range. The production-recommendation pattern from \S\ref{sec:results:pareto} (lower threshold for free savings) generalises in spirit; the cross-dataset specific recommendation is to calibrate the threshold per dataset rather than to assume \dsob{}'s default carries over.

\paragraphhead{Ablation.} The skill-ablation grid on \dsotel{} replicates \dsob{}'s pattern: removing the numeric-feature classifier collapses triage accuracy by $15.8$ absolute points (paired CI $[-0.207, -0.117]$, $p < 10^{-4}$), and no other single-primitive ablation moves Hit@K. The numeric-feature finding is therefore confirmed independently on a second dataset at the same significance level. The on-demand evidence-gathering loop's ablation cells, by contrast, produce \emph{exactly} zero per-window deltas on every metric: on \dsotel{} the loop fires on $25$ of $247$ windows ($10.1\%$, versus $\dsob{}$'s $17.8\%$), and on those $25$ windows the rerank step does not change the rank-one ranking. The ReAct loop's negative result from \dsob{} is sharper on \dsotel{}: not merely non-significant, but exactly zero.

\paragraphhead{Failure-mode catalog.} The tool-use catalog on \dsotel{} reports $25$ invocations of each tool. The pod-event tool returns an empty result on $22$ of $25$ ($88\%$) because the dataset's Kubernetes-event capture is sparser than \dsob{}'s; the other three tools succeed on all or all-but-two windows. The five-mode taxonomy (\S\ref{sec:results:failures}) thus produces a different per-tool distribution on \dsotel{}, but the same three never-firing modes (hallucinated, looping, budget-exhausted) remain at zero, confirming that the v1 controller's structural guarantees hold across datasets.

\subsection{Real Apache Jira tickets: the strongest external-validity test}
\label{sec:cross-app:wol}

The second transfer test is the paper's strongest external-validity claim. \dswol{} (\S\ref{sec:evaluation:datasets:wol}) replaces synthetic humanised tickets with $2{,}000$ real Apache Jira tickets from seven open-source projects, holds out two entire project families (Kafka and MariaDB Server) from training, and evaluates on the held-out projects' tickets paired with telemetry-style query text. This is a stricter holdout than \dsotel{}'s in two senses: the memory side is real production language rather than humanised synthetic prose, and the cross-project barrier means the test windows are drawn from projects the model has never seen during fine-tuning.

Table~\ref{tab:headline-wol} reports the agent's headline metrics.

\begin{table}[ht]
  \centering
  \caption{Headline metrics on \dswol{}, real Apache Jira tickets ($n = 304$, $n_{\text{eval}} = 55$).}
  \label{tab:headline-wol}
  \small
  \begin{tabular}{lr}
    \toprule
    Metric & Point \\
    \midrule
    Hit@1                      & $\mathbf{0.7818}$ \\
    Hit@5                      & $\mathbf{0.8364}$ \\
    MRR                        & $\mathbf{0.8045}$ \\
    Triage accuracy            & $\mathbf{0.9342}$ \\
    Novel recall / precision   & $0.0000 / 0.0000$ \\
    Pages-per-incident         & $\mathbf{1.000}$ \\
    Suppressions fired         & $18$ \\
    Distinct plan identifiers  & $1 / 6$ \\
    \bottomrule
  \end{tabular}
\end{table}

The headline metrics on \dswol{} are the highest of any dataset tested. Hit@1 of $0.7818$, Hit@5 of $0.8364$, MRR of $0.8045$, and triage accuracy of $0.9342$ all exceed \dsob{} and \dsotel{} by margins outside the latter datasets' confidence intervals. The cross-project, real-production-language test is in fact \emph{easier} for the agent than the two synthetic benchmarks, not harder---in part because the real-Jira ticket prose is more semantically discriminative per token (engineer-written prose is denser in fault-specific terminology than humanised synthetic prose), and in part because the dataset's evaluable subset of $55$ windows is a higher-precision sample under the time-ordered visibility rule.

\paragraphhead{The small evaluable subset.} The first table caveat is the $n_{\text{eval}} = 55$ figure. The held-out-projects design means most test windows have no time-visible gold ticket: their gold ticket exists in the corpus but is from a project family that the test split has explicitly excluded from visibility. We report the evaluation only on the $55$ windows that pass the visibility rule. The resulting confidence intervals are wider than \dsob{}'s ($\pm 5\%$ Hit@1 absolute on \dsob{} versus $\pm 9\%$ on \dswol{}) but the headline ordering remains stable: \dswol{} Hit@1 is the highest point estimate across the three datasets even at the lower CI bound.

\paragraphhead{Plan diversity collapses to one.} The controller emits exactly one distinct plan identifier across the $304$ \dswol{} test windows. The reason is structural rather than algorithmic: real Jira tickets do not carry the fault-window taxonomy (\texttt{pre\_fault\_baseline}, \texttt{recovery\_window}, \texttt{observation\_window}, \texttt{active\_fault}) that controls four of the six controller branches. Without a window-type signal, the controller defaults to the active-fault branch for every window, and the six-branch machinery is unexercised. This is not a failure---the agent runs end-to-end, the gates fire correctly within the active-fault branch, and the cross-window state layer suppresses repeated pages---but it is the most informative limitation of the controller's design: the per-window plan diversity claim is contingent on the dataset carrying a window-type taxonomy.

\paragraphhead{ReAct on \dswol{}.} Only one of the four on-demand evidence-gathering tools is structurally applicable on \dswol{}: the peer-incident retrieval tool, which reads the in-memory ticket corpus and does not depend on raw telemetry. The other three tools (Kubernetes events, extended trace, Prometheus snapshot) auto-drop because their required capability flags are absent in the dataset's bundles. The peer-incident tool's paired-delta on the $55$ evaluable windows is exactly zero: re-ranking does not change a single window's rank-one ranking. The cross-dataset ReAct picture is therefore consistent across all three datasets: a non-significant point-estimate lift on \dsob{}, an exact zero on \dsotel{}, an exact zero on \dswol{}. The framework is exercised; the value is below the noise floor at this sample size and operating point.

\paragraphhead{The verifier is structurally skipped on \dswol{}.} The dataset profile lists \dswol{} as \emph{known-harmful} for the language-model verifier, citing prior measurement that the verifier degrades Hit@5 on real Apache Jira tickets~\cite{tawos}. The capability gate (\S\ref{sec:approach:gating}) refuses to invoke the verifier on \dswol{} as a structural property, not a runtime decision; a unit-test assertion in the evaluation harness fires if the verifier ever appears in any \dswol{} trace's invocation list. This is the agent's structural-defence pattern: a measured negative result on one dataset (the verifier degrades retrieval on \dswol{}) becomes a configuration entry that the controller honours on future runs without re-measurement.

\subsection{What generalises across all three datasets}
\label{sec:cross-app:generalises}

Table~\ref{tab:generalises} summarises the cross-dataset pattern for each headline claim of the paper.

\begin{table*}[t]
  \centering
  \caption{Cross-dataset pattern. ``Generalises'' indicates the claim survives the dataset change with the same direction and similar magnitude; ``dataset-specific'' indicates a meaningful difference whose mechanism is identified in the corresponding section.}
  \label{tab:generalises}
  \small
  \begin{tabular}{llll}
    \toprule
    Claim                                                 & \dsob{}              & \dsotel{}             & \dswol{}              \\
    \midrule
    Cost savings vs.\ always-on cascade                   & $64.1\%$             & $65.3\%$              & verifier disabled by profile \\
    Pages-per-incident                                    & $1.000$              & $1.000$               & $1.000$               \\
    \texttt{triage\_numeric} essentiality ($p$)           & $<10^{-4}$           & $<10^{-4}$            & not applicable (no NUMERIC) \\
    Plan diversity (distinct plan IDs)                    & $4$ of $6$           & $4$ of $6$            & $1$ of $6$            \\
    Naive-baseline lift on Hit@5 (vs.\ BM25)              & $8.66\times$         & $1.34\times$          & $1.39\times$          \\
    \midrule
    \multicolumn{4}{l}{\emph{Dataset-specific:}} \\
    ReAct point-estimate lift on Hit@1                    & $+0.0151$ ($p=0.18$) & $0.0000$              & $0.0000$              \\
    Extended-trace tool harm                              & $\Delta=-0.018$, $p=0.002$ & dataset has trace flag; ablation flat & not applicable (no trace) \\
    False-novelty rate                                    & $40.6\%$             & $39.7\%$              & $\mathbf{0.7\%}$ \\
    \bottomrule
  \end{tabular}
\end{table*}

Three claims generalise without qualification across all three datasets and constitute the paper's strongest cross-dataset findings. First, the cost-savings claim holds at $64$ to $65$ percent against the always-on cascade counterfactual on both telemetry-rich datasets; on \dswol{} the verifier is structurally skipped, so the question does not arise. Second, the page-suppression mechanism converges to pages-per-incident $= 1.000$ on all three datasets, exercised by $1$ to $20$ suppression events per dataset; this is the most universal operational claim the paper makes. Third, the agent dominates a BM25-only naive baseline on Hit@5 by margins ranging from $1.34\times$ on \dsotel{} to $8.66\times$ on \dsob{}.

Two more claims generalise within their applicable scope. The numeric-feature classifier's essentiality for triage accuracy reaches $p < 10^{-4}$ on the two datasets that carry numeric features ($\dsob{}$, $\dsotel{}$); on \dswol{} the classifier is structurally absent because the dataset has no numeric-features capability. The controller's plan-diversity claim holds on the two telemetry-rich datasets that carry a fault-window taxonomy ($4$ of $6$ branches fire on each); on \dswol{} the controller collapses to one branch because the dataset has no window-type field.

\subsection{What does not generalise, and why}
\label{sec:cross-app:doesnt}

Three findings from \S\ref{sec:results} do not generalise.

\paragraphhead{The on-demand evidence-gathering lift.} The point-estimate Hit@1 lift from the ReAct loop on \dsob{} is $+0.0151$ ($p = 0.18$, non-significant). On \dsotel{} and \dswol{} the lift is exactly zero. The loop fires on a strict subset of windows on every dataset, but the rerank step's overlap-based scoring does not change rank-one rankings on \dsotel{} (where the pod-event tool returns empty on $88\%$ of fires and the peer corpus is sparse, $147$ tickets versus \dsob{}'s $347$) or on \dswol{} (where only the peer tool fires and the $55$-window subset offers too few candidate swaps to surface). The cross-dataset honest claim is that the framework is implemented and exercised; the empirical lift requires either a denser corpus or a more discriminative scoring head, both of which are future work.

\paragraphhead{The extended-trace tool's significant harm.} The $-0.018$ Hit@1 harm at $p = 0.002$ from the trace tool (\S\ref{sec:results:ablation}) is \dsob{}-specific. On \dsotel{} the same ablation cell produces zero per-window delta because the trace tool fires on too few windows ($25$ total) for an isolated-trace cell to surface a swap. On \dswol{} the tool does not fire at all because the trace-summary capability is absent. The harm finding is therefore dataset-specific in the sense that it is detectable only on the dataset where the tool is exercised most heavily, but the mechanism (cross-family token noise on call-path service names) is dataset-agnostic and the production recommendation (disable the tool by default; enable per scenario family) is portable.

\paragraphhead{The false-novelty rate.} The largest cross-dataset finding of the paper sits in the last row of Table~\ref{tab:generalises}. The false-novelty rate---the fraction of decisions in which the agent flags the window as novel despite a memory-side gold being present---is $40.6\%$ on \dsob{}, $39.7\%$ on \dsotel{}, and $0.7\%$ on \dswol{}. The two synthetic datasets agree to within one percentage point; the real-Jira dataset is sixty times lower. We had not predicted this and it is the cleanest cross-dataset discovery in the work.

\subsection{The false-novelty mechanism: a language-model-induced blind spot}
\label{sec:cross-app:falsenovel}

The mechanism is the structural skip of the language-model verifier on \dswol{} (\S\ref{sec:cross-app:wol}). The novelty composition layer (\S\ref{sec:approach:primitives}) OR-combines three independent novelty signals: the verifier's emitted novelty flag, a maximum-retrieval-confidence threshold, and a learned per-window classifier conditioned on coarse window metadata. The verifier's signal is the one that does not exist on \dswol{} because the verifier itself does not run. On \dsob{} and \dsotel{} the verifier runs and contributes its signal; on the windows where the verifier-induced novelty flag is over-eager (the verifier sees only the top-$k$ retrieved candidates and judges that none of them match the hypothesised root cause), the OR composition fires and the agent emits a false-novelty decision.

We test the mechanism directly. On \dsob{}, the windows that fire the false-novelty bucket overlap by $94\%$ with the windows on which the verifier emits its own novelty flag, and by under $10\%$ with windows where the other two signals would fire on their own. The verifier is the source of the false-novelty signal on the two synthetic datasets, not the OR-composition's failure to reject the verifier's input.

The implication is the paper's clearest design recommendation for future work. The verifier's novelty-flag emission is a tunable: it can be replaced by a more conservative criterion (e.g., requiring two of the three novelty signals to agree rather than one), or it can be disabled selectively on operating points where the verifier is not known-helpful for novelty. The latter is a configuration change; the former is a one-line edit to the novelty composer's OR-disjunction. Either change would, by the same mechanism that produces \dswol{}'s $0.7\%$ false-novelty rate, collapse \dsob{} and \dsotel{} false-novelty by a similar factor. The full evaluation of the change is left to a follow-up. The finding itself is what \S\ref{sec:cross-app}'s replication makes possible: a single-dataset paper would not have surfaced it.

\subsection{Summary of the cross-dataset evidence}
\label{sec:cross-app:summary}

The agent's design transfers across applications and across the synthetic-to-real boundary without per-dataset tuning. The strongest cross-dataset claims are the universal page-suppression ratio, the cost-savings rate at $64$ to $65$ percent, and the three-dataset paired-delta significance of the numeric-feature classifier on triage accuracy. The Hit@K numbers on real Apache Jira tickets exceed both synthetic benchmarks, in part because real engineer-written ticket prose carries more discriminative signal per token than humanised synthetic prose. The on-demand evidence-gathering loop's Hit@K contribution is not statistically distinguishable from zero on any of the three datasets, and the cross-dataset replication of the negative finding is what makes the result reportable rather than mistakable for a within-dataset accident. The false-novelty rate's sixty-fold collapse on \dswol{} relative to the synthetic datasets is the paper's cleanest mechanistic discovery and points to a tractable design fix for the agent's largest blind spot.
